% ECONOMICS_PROBLEM: {"id": "E31-467", "title": "Supply-versus-demand inflation", "jel_code": "E31", "source_id": "OP-0241"}
% !TeX program = xelatex
\documentclass[11pt,a4paper]{article}
\usepackage[margin=23mm]{geometry}
\usepackage{fontspec}
\setmainfont{Latin Modern Roman}
\usepackage{amsmath,amssymb,mathtools}
\usepackage{xcolor,enumitem,booktabs,longtable,array,xurl,hyperref}
\definecolor{muted}{HTML}{53636C}
\hypersetup{colorlinks=true,urlcolor=blue,linkcolor=blue}
\setlength{\parindent}{0pt}
\setlength{\parskip}{5pt}
\newcommand{\R}{\mathbb R}
\newcommand{\E}{\mathbb E}
\newcommand{\N}{\mathbb N}
\newcommand{\ind}{\mathbf 1}
\newcommand{\Var}{\operatorname{Var}}
\newcommand{\Cov}{\operatorname{Cov}}
\newcommand{\supp}{\operatorname{supp}}
\DeclareMathOperator*{\argmin}{arg\,min}
\DeclareMathOperator*{\argmax}{arg\,max}
\newcommand{\entrymeta}[3]{{\sffamily\small\color{muted}\textbf{\texttt{#1}}\quad|\quad #2\par #3\par}}
\newcommand{\Problemref}[1]{\texttt{#1}}
\newcommand{\JELref}[2]{\texttt{#2} (JEL \texttt{#1})}
\begin{document}
\section*{Supply-versus-demand inflation}
\textbf{Problem ID:} \texttt{E31-467}\quad \textbf{JEL:} \texttt{E31}\par
% BEGIN ECONOMICS STATEMENT
\label{problem:OP-0241}
{\sffamily\small\color{muted}Original subject group: Inflation and monetary policy\par}
\label{definitions:problem:30007522}
\textbf{Audit classification:} Background; proposed extension. An inflation-driver decomposition is already provided; robustness and Shapley allocation here are editorial extensions. Verification concerns the cited version, not first priority or proof that this exact editorial specification remains unsolved on October 8, 2026.
\subsection*{Definitions and precise research target}
\paragraph{Editorial mathematical specification.} Fix a multicountry, multisector economy with observed sectoral prices, wages, output, expenditures, and input–output links. Its structural model maps a finite sequence of shocks \(e\) and parameter \(\theta\) into aggregate inflation \(\pi_t(e;\theta)\). Partition shocks into supply, aggregate demand, and sectoral demand reallocation. State the instruments or exclusion restrictions distinguishing these classes; price–quantity comovement alone must not be assumed to establish causal classification. Let \(e^S\) retain only shocks in subset \(S\) of the three classes, with a common initial state and policy rule.
To allocate nonlinear interactions, define a Shapley contribution for class \(k\):
\[
C_{k,t}=\sum_{S\subseteq K\setminus\{k\}}
\frac{|S|!(|K|-|S|-1)!}{|K|!}
\big[\pi_t(e^{S\cup\{k\}};\theta)-\pi_t(e^S;\theta)\big],
\]
where \(K\) is the three-class set. Identify or bound these contributions over a declared inflation episode, including uncertainty in shocks, elasticities, and network links. Test whether decompositions remain stable under plausible alternative identification schemes and inflation aggregation weights. Counterfactual policy responses must be declared as fixed rules or endogenous responses. The target is an explicitly convention-dependent causal decomposition, not a unique model-free attribution of each observed price increase to supply or demand.

Define aggregate inflation by a fixed index \(P_t>0\), \(\pi_t=\log P_t-\log P_{t-1}\), or disclose a different index and its weights. A supply innovation changes a production or factor-endowment primitive; an aggregate-demand innovation changes total desired expenditure; a reallocation innovation changes expenditure shares with an explicit normalization. Retaining a shock class means retaining its entire predeclared time path, setting the others to their zero-innovation values, and solving equilibrium anew. With \(v_t(S)=\pi_t(e^S;\theta)\), the displayed formula implies \(\sum_kC_{k,t}=v_t(K)-v_t(\varnothing)\). Contributions may be negative and need not sum to observed inflation unless the baseline and all omitted disturbances are accounted for. The causal target is the set of contributions over structurally compatible parameters and shock paths, not the Shapley arithmetic itself.
\subsection*{Variants and assumption changes}
The following are \emph{editorial variants}, not additional published open conjectures.
\begin{enumerate}
\item \emph{Linear benchmark.} Impose an affine equilibrium map in the shock paths. Then cross-class interactions vanish and Shapley contributions equal ordinary additive class effects; investigate identification of the shocks rather than treating this accounting identity as open.
\item \emph{Policy-response convention.} Compare decompositions under the same fixed monetary rule with decompositions under reoptimized policy in every counterfactual. The latter includes an induced policy effect and answers a different intervention question.
\end{enumerate}
\subsection*{References and source audit}
\begin{itemize}
\item Julian di Giovanni; \c{S}ebnem Kalemli-\"{O}zcan; Alvaro Silva; Muhammed A. Y\i{}ld\i{}r\i{}m. \emph{Pandemic-Era Inflation Drivers and Global Spillovers}. 2023. Locator: Abstract and Section1, printed pp.1--3; model-based supply/demand decomposition. Primary text checked. \url{https://www.newyorkfed.org/medialibrary/media/research/staff_reports/sr1080.pdf}.
\end{itemize}
An inflation-driver decomposition is already provided; robustness and Shapley allocation here are editorial extensions.
\begin{enumerate}\raggedright
\item\hypertarget{definitions:source:30007522:1}{} Julian di Giovanni; \c{S}ebnem Kalemli-\"{O}zcan; Alvaro Silva; Muhammed A. Y\i{}ld\i{}r\i{}m. 2023. \emph{Pandemic-Era Inflation Drivers and Global Spillovers}. Abstract and Section1, printed pp.1--3; model-based supply/demand decomposition.
\url{https://www.newyorkfed.org/medialibrary/media/research/staff_reports/sr1080.pdf}.
DOI: \url{https://doi.org/10.59576/sr.1080}.
\end{enumerate}

\subsection*{Common conventions: Source mathematical conventions}

These conventions apply to each entry unless explicitly replaced.

\textbf{Models and identification.} A primitive model \(m\) specifies all probability laws, feasible choices, information, equilibrium and measurement rules used by its target. The admissible class \(\mathcal M\) is fixed independently of outcomes. If \(P_O(m)\) is its observable probability law and \(\tau(m)\) the target, its identified set at observable law \(P\) is
\[
 \mathcal I_\tau(P)=\{\tau(m):m\in\mathcal M,\ P_O(m)=P\}.
\]
Point identification means that this set is a singleton. An empty set signals incompatibility of data and assumptions. Coordinate bounds need not describe the joint set. Bounds are sharp only if every retained target is attainable or arbitrarily approachable by compatible models. Finite bounds require explicit restrictions on unobserved quantities; unrestricted confounding, hidden wealth, extrapolation or arbitrary equilibrium selection can yield uninformative sets.

\textbf{Causal interventions.} A potential outcome \(Y_i(p)\) is the outcome under a complete policy regime \(p\), including timing, financing, information and market spillovers. The target population and units are fixed before treatment. With interference, use the complete assignment vector or a justified exposure mapping; individual no-interference notation alone cannot identify an economy-wide intervention. A difference of mean potential outcomes does not require the cross-world joint distribution. Conditioning on post-treatment migration, survival, enrollment or employment changes the estimand and needs separate justification.

\textbf{Statistical statements.} An estimator, confidence set, forecast or test specifies a sampling law, model class and loss/coverage criterion. Uniform \((1-\alpha)\)-coverage means \(\inf_{m\in\mathcal M}\Pr_m\{\tau(m)\in C_n\}\geq1-\alpha\), or an explicitly asymptotic version. The whole parameter space trivially has coverage, so informativeness requires a stated width, risk or power objective. A forecasting improvement is relative to a named benchmark, information set and evaluation population.

\textbf{Equilibrium and optimization.} An equilibrium comprises feasible plans, optimality, beliefs where needed, budget constraints and all market/resource clearing. The primitives or a stated benchmark define each correspondence \(\mathcal E(p;m)\). Nonemptiness is assumed or proved, never inferred just from compact policy sets. A selected-equilibrium target uses a declared selection rule; a robust target quantifies over all permitted equilibria. Use suprema or infima unless attainment is established. An \(\varepsilon\)-optimal policy is feasible and within \(\varepsilon\) of the finite optimum value. Report an empty feasible set separately.

\textbf{Welfare and accounting.} Normative utility, interpersonal normalization, social weights, numeraire, discounting and the use of public receipts are primitives. Behavioral utility need not coincide with the welfare criterion. Household welfare includes ownership income and rents once; adding profits again to compensating variation generally double-counts them. A monetary equivalent is defined by an explicit transfer and price convention, with monotonicity and range conditions. Average effects, mechanism contributions, transfers and welfare losses are different targets. Infinite sums require convergence and dynamic feasible plans require terminal or no-Ponzi conditions.

\textbf{Source versions.} A working-paper date, revision date and journal publication year are distinguished. A DOI for a journal version is not automatically the DOI of an earlier manuscript. Locators refer to the stated version and printed pages where specified. A metadata-only check does not verify a claim about a full-text conclusion. Every entry's source subsection narrows the provenance claim accordingly.
% END ECONOMICS STATEMENT
\end{document}
